\documentclass[12pt,a4paper]{article}
\usepackage[T1]{fontenc}
\usepackage[utf8]{inputenc}
\usepackage[ngerman]{babel}
\usepackage{geometry}
\usepackage{enumitem}
\usepackage{hyperref}
\geometry{margin=2.5cm}
\setlist{nosep}
\hypersetup{hidelinks}
\emergencystretch=3em

\title{Exposé zur Bachelorarbeit\\[0.5em]
\large Robustheit von TabPFN-3.5 gegenüber irrelevanten Merkmalen:
Eine kontrollierte Untersuchung im Vergleich zu LightGBM}
\author{}
\date{Vorläufiger Entwurf}

\begin{document}
\maketitle

\section{Problemstellung und Motivation}

Tabellarische Klassifikationsaufgaben unterscheiden sich in Umfang,
Merkmalstypen, Klassenverteilung und Datenqualität. Für solche Aufgaben
werden häufig baumbasierte Verfahren eingesetzt. Grinsztajn et al.\
untersuchen in einem Benchmark mit 45 kuratierten Datensätzen verschiedene
baumbasierte und neuronale Verfahren. Sie berichten unter anderem, dass das
Hinzufügen unkorrelierter irrelevanter Merkmale die Leistung der untersuchten
neuronalen Modelle beeinträchtigen und den Abstand zu den Baumverfahren
vergrößern kann \cite{grinsztajn_why_2022}. Der Befund gilt für die in dieser
Studie getesteten Modelle und Versuchsbedingungen, nicht automatisch für
jede neuronale Architektur oder jede spätere Modellgeneration.

TabPFN verfolgt einen anderen Ansatz als ein pro Datensatz neu trainiertes
neuronales Netz: Ein Transformer wird vorab auf synthetisch erzeugten
Klassifikationsaufgaben trainiert und verarbeitet beim Einsatz die
beschrifteten Trainingsbeispiele als Kontext. Das ursprüngliche Paper
evaluiert TabPFN vor allem auf kleinen Klassifikationsdatensätzen und
untersucht zusätzlich den Einfluss uninformativer Merkmale. Für die damalige
Modellversion berichten die Autoren eine geringere Robustheit gegenüber
solchen Merkmalen als bei LightGBM \cite{hollmann_tabpfn_2023}. Dieser
Befund bietet einen konkreten Ausgangspunkt für die vorliegende Arbeit.

Der technische Bericht zu TabPFN-3.5 beschreibt unter anderem Anpassungen
des synthetischen Priors für Datensätze mit größeren Merkmalszahlen und
berichtet Benchmark-Ergebnisse für hochdimensionale Daten
\cite{jager_tabpfn-35_2026}. Der Bericht ist eine Evaluation des
Modellentwicklerteams und wird in der Arbeit entsprechend als solcher
eingeordnet. Eine größere unterstützte Merkmalszahl beantwortet für sich
genommen nicht, wie empfindlich TabPFN-3.5 reagiert, wenn gezielt zusätzliche
Merkmale ohne Zusammenhang mit dem Ziel eingeführt werden. Im geprüften
Bericht wird kein isoliertes Experiment mit kontrolliert hinzugefügten
unabhängigen Störmerkmalen berichtet. Ob neuere unabhängige oder
anwendungsbezogene Studien diese konkrete Frage bereits untersuchen, wird
vor einer abschließenden Behauptung einer Forschungslücke systematisch
nachgeprüft.

Damit ergibt sich eine klar begrenzte praktische Frage: Verhält sich
TabPFN-3.5 bei kontrollierter Zugabe irrelevanter Merkmale anders als
LightGBM, und verändert sich die Vorhersageleistung der Modelle mit
zunehmendem Störmerkmal-Anteil? Die Untersuchung soll keinen allgemeinen
Sieger zwischen Foundation Models und Baumverfahren bestimmen. Sie isoliert
stattdessen einen Datenfaktor, der in früheren Arbeiten als relevant
beschrieben wurde.

\section{Zielsetzung und möglicher Beitrag}

Ziel der Arbeit ist es, ausgewählte methodische Änderungen von TabPFN bis
Version 3.5 und die Aussagekraft relevanter Evaluationen knapp und kritisch
einzuordnen. Darauf aufbauend wird in einem kontrollierten
Klassifikationsexperiment untersucht, wie sich die Leistung von TabPFN-3.5
und LightGBM verändert, wenn unabhängig erzeugte irrelevante Merkmale
schrittweise zu den Eingabedaten hinzugefügt werden.

Der praktische Beitrag besteht in einer reproduzierbaren Untersuchung mit
mehreren geeigneten Datensätzen, festgelegten Störmerkmal-Stufen,
wiederholten Datenaufteilungen und einer klar benannten Baum-Baseline.
Gemessen wird die Veränderung gegenüber der jeweiligen Bedingung ohne
zusätzliche Störmerkmale. Dadurch steht die Empfindlichkeit gegenüber dem
kontrollierten Eingriff im Vordergrund und nicht allein ein Rangvergleich
der Modelle auf unterschiedlichen Datensätzen. Die Arbeit entwickelt keinen
neuen Algorithmus und beansprucht nicht, Ergebnisse für alle baumbasierten
Verfahren zu liefern. Ob die Untersuchung eine Forschungslücke schließt,
hängt vom noch ausstehenden Abgleich mit der neueren Literatur ab; bis dahin
wird sie als empirische Prüfung und mögliche Erweiterung früherer Befunde
bezeichnet.

\section{Forschungsfragen}

\subsection*{Übergeordnete Forschungsfrage}

\textbf{Wie robust ist TabPFN-3.5 gegenüber schrittweise hinzugefügten
unabhängigen, irrelevanten Merkmalen im Vergleich zu LightGBM, und wie ist
dieses Ergebnis anhand der methodischen Entwicklung und bisherigen
Evaluationen von TabPFN einzuordnen?}

Die übergeordnete Frage verbindet Literatur- und Experimentteil. Die
Literaturanalyse klärt, wie das Modell und die Fragestellung einzuordnen
sind; der experimentelle Teil beantwortet den konkreten Leistungsvergleich.
Die übergeordnete Frage verspricht keine experimentelle Prüfung sämtlicher
TabPFN-Entwicklungsschritte.

\subsection*{FF1: Literaturanalyse}

\textbf{Welche relevanten methodischen Änderungen und adressierten
Herausforderungen tabellarischer Daten beschreiben ausgewählte
TabPFN-Arbeiten, und welche Schlussfolgerungen erlauben ihre Evaluationen?}

\subsection*{FF2: Experimenteller Vergleich}

\textbf{Wie verändert sich die Klassifikationsleistung von TabPFN-3.5 und
LightGBM bei schrittweiser Zugabe unabhängiger, irrelevanter Merkmale?}

FF2 wird anhand der Veränderung der vorab festgelegten Leistungsmetriken
über Störmerkmal-Stufen, Datensätze und Wiederholungen beantwortet. Aus dem
Vergleich mit LightGBM werden keine Aussagen über andere Baumverfahren
abgeleitet.

\section{Inhaltliche Abgrenzung}

Die Untersuchung konzentriert sich auf überwachte tabellarische
Klassifikation und auf den Einfluss künstlich hinzugefügter Merkmale, die
unabhängig von den ursprünglichen Eingaben und der Zielvariable erzeugt
werden. Der Modellvergleich umfasst TabPFN-3.5 und LightGBM. Die genaue
Version, Paketkonfiguration und Verfügbarkeit von TabPFN-3.5 werden vor
Beginn des Hauptversuchs technisch und lizenzrechtlich geprüft.

Nicht Gegenstand der Arbeit sind Explainable AI, SHAP oder LIME,
Kausalaussagen, Regression, ein umfassender Benchmark vieler Modellfamilien
oder eine vollständige Versionsgeschichte von TabPFN. Laufzeit wird nicht
als zentrale Forschungsfrage behandelt; eine ergänzende Laufzeitmessung
kann nur aufgenommen werden, wenn sie unter dokumentierten und fairen
Bedingungen realisierbar ist. Vorgängermodelle werden ausschließlich dort
herangezogen, wo sie für die Einordnung der aktuellen Version und früherer
Robustheitsbefunde nötig sind.

\section{Roter Faden und inhaltlicher Aufbau}

Der rote Faden führt vom Problem zu einer gezielten empirischen Prüfung:

\begin{quote}
Tabellarische Daten enthalten Merkmale mit unterschiedlichem Informationswert
$\longrightarrow$ frühere Benchmarks und TabPFN-Arbeiten berichten einen
Einfluss irrelevanter Merkmale auf die Leistung bestimmter Modelle
$\longrightarrow$ TabPFN-3.5 erweitert laut technischem Bericht den Umgang
mit größeren Merkmalszahlen, doch daraus folgt nicht automatisch Robustheit
gegenüber gezielt hinzugefügten Störmerkmalen
$\longrightarrow$ die Arbeit variiert diese Merkmale kontrolliert und
vergleicht die Veränderung der Klassifikationsleistung von TabPFN-3.5 und
LightGBM.
\end{quote}

Grundlagen zu tabellarischem Deep Learning, Prior-Data Fitted Networks und
dem synthetischen Prior von TabPFN dienen dazu, die Modellidee und den
Bezug der Forschungsfrage zu erklären. Sie werden nicht zu einem breiten
allgemeinen Abriss über KI oder LLMs ausgebaut. Die Entwicklung früherer
Versionen ist selektiver Kontext statt eines eigenen Kapitels zur
Versionsgeschichte. Der Schwerpunkt und die Eigenleistung liegen im
kontrollierten Experiment und dessen kritischer Auswertung.

\section{Vorgehen der Literaturanalyse}

Die Literaturanalyse wird als fokussierte, nachvollziehbare Recherche
angelegt, nicht als Anspruch auf ein vollständiges systematisches Review
aller Tabular Foundation Models. Im Zentrum stehen:
\begin{itemize}
  \item das ursprüngliche TabPFN-Paper und die darin berichteten
        Evaluationen, insbesondere der Befund zu uninformativen Merkmalen;
  \item der Benchmark von Grinsztajn et al.\ als Kontext für den Vergleich
        neuronaler Modelle und Baumverfahren;
  \item der TabPFN-3.5-Bericht mit Fokus auf die für Feature-Zahl,
        synthetischen Prior und Evaluation relevanten Aussagen;
  \item neuere Arbeiten, die gezielt TabPFN-3/3.5, irrelevante Merkmale,
        Feature-Robustheit oder kontrollierte Feature-Interventionen
        untersuchen.
\end{itemize}

Gesucht wird unter anderem in Google Scholar, arXiv und den
Literaturverzeichnissen einschlägiger Konferenzbeiträge. Vorläufige
Suchkombinationen kombinieren \texttt{TabPFN} mit \texttt{irrelevant
features}, \texttt{robustness}, \texttt{noise features} oder
\texttt{feature selection}; zusätzlich wird nach \texttt{tabular foundation
model} zusammen mit \texttt{robustness} und \texttt{irrelevant features}
gesucht. Suchdatum, Suchanfragen, relevante Treffer und Ausschlüsse
werden festgehalten. Bei der Synthese werden Modellversion, Datensatz,
Versuchsaufbau, Metriken und Grenzen der jeweiligen Quelle gegenübergestellt.
Entwicklerberichte werden als solche kenntlich gemacht; berichtete
Benchmarkresultate werden nicht ungeprüft als unabhängige Evidenz
interpretiert.

\section{Vorläufiger experimenteller Vergleich}

\subsection*{Datensätze und Versuchsbedingungen}

Vorgesehen sind zunächst etwa vier öffentlich verfügbare
numerische Klassifikationsdatensätze. Die Beschränkung auf numerische
Eingaben soll zusätzliche Effekte aus unterschiedlichen Verfahren zur
Kodierung kategorialer Merkmale vermeiden. Die Datensätze werden anhand
vorab festgelegter Kriterien
ausgewählt: Eignung für TabPFN-3.5, Anzahl der Beobachtungen und Merkmale,
Klassenverteilung, Datenqualität und verfügbare Nutzungsrechte. Die
Datensatzauswahl und Ausschlüsse werden dokumentiert. Das genaue Set wird
erst nach technischer Prüfung festgeschrieben.

Für jeden Datensatz werden identische, stratifizierte Datenaufteilungen für
beide Modelle verwendet. Auf die Eingabemerkmale werden in kontrollierten
Stufen zusätzliche numerische Störmerkmale gelegt, die unabhängig von den
ursprünglichen Merkmalen, voneinander und der Zielvariable erzeugt werden.
Als erste Versuchsplanung sind unabhängige standardnormalverteilte
Störmerkmale und Stufen von 0, 0{,}5, 1 und 2 zusätzlichen Merkmalen pro
ursprünglicher Eingabespalte vorgesehen; nicht-ganzzahlige Stufenzahlen
werden auf ganze Spalten gerundet. Die genaue Definition der ursprünglichen
Eingabespalten und die Machbarkeit der höchsten Stufe werden im Pilotlauf
geprüft. Dieselben erzeugten Merkmale, Seeds und Splits werden für beide
Modelle verwendet.

Die Auswertung wird über mehrere wiederholte Seeds durchgeführt. Eine
unveränderte Bedingung ohne zusätzliche Merkmale dient pro Datensatz und
Modell als Referenz. Die experimentellen Bedingungen unterscheiden sich
innerhalb eines Vergleichs nur durch die Zahl der Störmerkmale.

\subsection*{Modelle, Metriken und Auswertung}

TabPFN-3.5 wird mit dokumentierter Paket- und Modellversion sowie
festgehaltenen Inferenzoptionen ausgewertet. LightGBM wird als einzelnes
baumbasiertes Gradient-Boosting-Modell eingesetzt \cite{ke2017lightgbm}.
Um Robustheit statt zusätzlicher Tuningeffekte zu vergleichen, werden
Modellkonfigurationen vor der Auswertung der Störmerkmal-Stufen festgelegt.
Für LightGBM kann die Konfiguration auf einer Validierung der unveränderten
Trainingsdaten bestimmt und anschließend über alle Stufen eingefroren
werden; TabPFN wird mit seiner dokumentierten Standardkonfiguration
verwendet. Diese asymmetrische Trainingsweise wird als Teil des Designs
offengelegt: Verglichen wird die Veränderung unter demselben Eingriff, nicht
der Aufwand einer gleich großen Hyperparametersuche.

Als primäre Kennzahl ist zunächst Balanced Accuracy vorgesehen, um
Klassifikationsleistung über Datensätze mit unterschiedlichen
Klassenverteilungen vergleichbar auszuwerten. Log-Loss kann ergänzend
berichtet werden, sofern die Wahrscheinlichkeitsausgaben technisch
vergleichbar und ohne zusätzliche, unterschiedlich abgestimmte Kalibrierung
auswertbar sind. Für jedes Modell und jeden Datensatz wird insbesondere die
Änderung gegenüber der Bedingung ohne Störmerkmale dargestellt. Ergebnisse
pro Datensatz, Streuung über Wiederholungen und aggregierte Ergebnisse
werden getrennt ausgewiesen. Unsicherheit wird angemessen berichtet; eine
aggregierte Rangliste allein ist nicht ausreichend.

\subsection*{Machbarkeit und Grenzen}

Vor dem Hauptversuch wird ein Pilot mit mindestens einem Datensatz
durchgeführt. Er prüft Installation und Lizenzbedingungen von TabPFN-3.5,
Laufzeit, Speicherbedarf, Datenformat und die geplante Vorverarbeitung.
Datensätze, Störmerkmal-Stufen und Wiederholungszahlen werden nur dann
festgelegt, wenn sie mit der verfügbaren Hardware realisierbar sind.

Unabhängig generierte Merkmale sind in endlichen Stichproben nicht zwingend
empirisch unkorreliert. Wiederholte Seeds und die kontrollierte
Erzeugungsprozedur helfen, zufällige Stichprobeneffekte sichtbar zu machen,
beseitigen sie aber nicht vollständig. Die Resultate gelten außerdem für
die ausgewählten Datensätze, die festgelegte Datenpipeline und genau die
beiden untersuchten Modelle. Sie erlauben keine direkte Verallgemeinerung
auf alle Datendomänen oder alle baumbasierten Verfahren.

\section{Vorläufige Gliederung}

\begin{enumerate}
  \item \textbf{Einleitung:} Problemstellung, Motivation, Zielsetzung,
        Forschungsfragen, Abgrenzung und Aufbau.
  \item \textbf{Grundlagen:} Tabellarische Klassifikation; TabPFN und
        Prior-Data Fitted Networks; synthetische Priors; Gradient Boosting
        und LightGBM; irrelevante Merkmale und Robustheitsbegriffe.
  \item \textbf{Stand der Forschung und Literaturanalyse:} Ausgewählte
        methodische Änderungen bis TabPFN-3.5; Benchmarkbefunde und deren
        Grenzen; frühere Experimente mit irrelevanten Merkmalen;
        Herleitung des Untersuchungsbedarfs.
  \item \textbf{Methodik:} Forschungsdesign; Datensatzauswahl;
        Störmerkmal-Erzeugung; Modelle und Konfigurationen; Splits,
        Wiederholungen, Metriken und Auswertung.
  \item \textbf{Ergebnisse:} Ergebnisse ohne Störmerkmale;
        Leistungsänderung mit zunehmender Störmerkmal-Zahl; Vergleich
        zwischen Modellen und Datensätzen.
  \item \textbf{Diskussion:} Beantwortung der Forschungsfragen;
        Einordnung in die Literatur; Bedeutung und Limitationen.
  \item \textbf{Fazit und Ausblick.}
\end{enumerate}

\bibliographystyle{plain}
\bibliography{references}

\end{document}